\chapter{Le cadre d'Almgren--Chriss}\label{ch:mx:the-almgren-chriss-framework}

Un trader doit vendre un million d'actions d'ici la clôture. Tout vendre immédiatement coûte l'impact~; les vendre régulièrement fait porter le risque de prix toute la journée. Entre
les deux s'étend une courbe de calendriers d'exécution, et choisir un point sur cette courbe revient à choisir une aversion au risque. Ce chapitre pose le problème comme l'ont fait Almgren et Chriss,
le résout sous forme fermée, trace la frontière, puis fait tourner les calendriers sur le marché simulé pour voir ce que le modèle décrit correctement et ce qu'il manque.

\section{Ordres parents, ordres enfants et une trajectoire}

\begin{definition}[Ordre parent, ordre enfant]\label{def:mx:the-almgren-chriss-framework:parent}
Un \emph{ordre parent}\index{ordre parent} est une instruction de négocier une quantité sur une période, confiée à un trader ou à un algorithme~; les
\emph{ordres enfants}\index{ordre enfant} sont les ordres effectivement envoyés aux plateformes pour l'exécuter.
\end{definition}

\begin{definition}[Trajectoire d'exécution]\label{def:mx:the-almgren-chriss-framework:trajectory}
Une \emph{trajectoire d'exécution}\index{trajectoire d'exécution} est la quantité $x(t)$ restant à négocier à chaque instant $t$ de la fenêtre d'exécution $[0,T]$, depuis
$x(0)=X$ jusqu'à $x(T)=0$~; son taux de négociation est $v(t)=-\dot x(t)$.
\end{definition}

L'\omterm{def:mx:the-almgren-chriss-framework:parent}{ordre parent} est un métaordre (chapitre~11) vu de l'intérieur. Un calendrier choisit la trajectoire~; le placement des \omterm{def:mx:the-almgren-chriss-framework:parent}{ordres enfants} (chapitre~17) et le routage
(chapitre~18) décident comment chaque tranche est négociée. Bertsimas et Lo (1998) ont posé le choix comme un problème de programmation dynamique minimisant le coût espéré~; Almgren et
Chriss (2001) ont ajouté le risque.

\section{Coûts~: impact et risque}

\begin{definition}[Modèle d'Almgren--Chriss]\label{def:mx:the-almgren-chriss-framework:ac}
Dans le \emph{modèle d'Almgren--Chriss}\index{modèle d'Almgren--Chriss}, le prix évolue sous l'effet d'un impact permanent $\gamma$ par action négociée et d'une marche aléatoire de
volatilité $\sigma$, et chaque négociation supporte un impact temporaire $\eta v$ par action au taux $v$. Vendre $X$ selon $x(t)$ coûte, par rapport au prix d'arrivée,
\[
\E[C]=\tfrac12\gamma X^2+\eta\int_0^Tv(t)^2\,dt,\qquad \operatorname{Var}[C]=\sigma^2\int_0^Tx(t)^2\,dt.
\]
\end{definition}

\begin{definition}[Risque d'exécution]\label{def:mx:the-almgren-chriss-framework:risk}
Le \emph{risque d'exécution}\index{risque d'exécution} est la variance (ou l'écart-type) du coût d'une exécution qui provient de la variation du prix pendant que la
position est toujours détenue~: $\sigma^2\int x^2\,dt$ ici.
\end{definition}

Le terme permanent ne dépend pas du calendrier. Le terme temporaire est minimal lorsqu'on négocie aussi lentement que possible, uniformément sur toute la fenêtre~; le
risque est minimal lorsqu'on négocie d'un coup. Un trader averse au risque minimise $\E[C]+\lambda\operatorname{Var}[C]$ (moyenne--variance, livre~7, chapitre~25),
avec l'aversion au risque $\lambda$ exprimée en inverse de monnaie.

\section{La trajectoire optimale sous forme fermée}

\begin{proposition}[Trajectoire d'Almgren--Chriss]\label{prop:mx:the-almgren-chriss-framework:sinh}
La trajectoire qui minimise $\E[C]+\lambda\operatorname{Var}[C]$ est
\[
x(t)=X\,\frac{\sinh\bigl(\kappa(T-t)\bigr)}{\sinh(\kappa T)},\qquad \kappa=\sqrt{\frac{\lambda\sigma^2}{\eta}}.
\]
\end{proposition}

\begin{proof}
L'objectif vaut $\int_0^T(\eta\dot x^2+\lambda\sigma^2x^2)\,dt$ à une constante près. Son équation d'Euler--Lagrange est $\eta\ddot x=\lambda\sigma^2x$, dont
les solutions sont des combinaisons de $e^{\pm\kappa t}$~; les conditions aux limites $x(0)=X$, $x(T)=0$ sélectionnent le rapport de sinus hyperboliques. (C'est aussi la voie de la fonction valeur
du livre~4, chapitre~9, pour un problème linéaire-quadratique.)
\end{proof}

\begin{definition}[Urgence]\label{def:mx:the-almgren-chriss-framework:urgency}
L'\emph{urgence}\index{urgence} d'une exécution est $\kappa$, ou le nombre sans dimension $\kappa T$~: nulle pour une droite (pondérée par le temps), grande pour une
trajectoire qui vend l'essentiel de l'ordre tôt. Son inverse est l'échelle de temps naturelle de l'exécution.
\end{definition}

Sur $n$ intervalles de longueur $\tau$, l'optimum discret a la même forme, avec $\tilde\kappa$ donné par $\cosh(\tilde\kappa\tau)=1+\lambda\sigma^2\tau^2/(2\eta)$.
La fonction \texttt{firm.acexec.discrete} l'implémente avec le coût et la variance de tout calendrier discret.

\omcode{code/firm/acexec/firm_acexec.py}{40}{57}{La trajectoire d'Almgren--Chriss discrète, ainsi que le coût espéré et la variance d'un calendrier.}\label{lst:mx:the-almgren-chriss-framework:ac}

Le cas d'étude du chapitre~: une action cotée 50~USD, qui s'échange à raison de 10 millions d'actions par jour avec une volatilité quotidienne de 2~\%, soit $\sigma=1$ dollar par action et par racine carrée de jour~;
vendre un million d'actions d'ici la clôture. L'impact temporaire est linéarisé à partir de la loi en racine carrée du chapitre~11 pour une participation de 10~\%,
$\eta=0.8\times0.02\times\sqrt{0.1}\times50/(0.1\times10^7)=2.53\times10^{-7}$ dollars par action et par jour~; pas d'impact permanent. Avec $\lambda=10^{-6}$ par
dollar, $\kappa=1.99$ par jour. Le calendrier optimal sur des intervalles d'une minute a un coût espéré de 300\,000~USD (30 cents par action, 60 points de base), avec un
écart-type de 470\,000~USD, un équivalent certain $\E+\lambda\operatorname{Var}$ de 521\,000~USD, et vend 95~\% de l'ordre en 5.9 des 6.5 heures
(la \cref{fig:mx:the-almgren-chriss-framework:frontier}, à gauche).

\section{La frontière d'exécution efficiente}

\begin{definition}[Frontière d'exécution efficiente]\label{def:mx:the-almgren-chriss-framework:frontier}
La \emph{frontière d'exécution efficiente}\index{frontière d'exécution efficiente} est l'ensemble des couples (risque, coût espéré) des trajectoires optimales lorsque l'aversion
au risque varie~: aucun calendrier n'a à la fois un risque et un coût espéré plus faibles qu'un point de la frontière.
\end{definition}

La droite est l'extrémité de la frontière à plus faible coût espéré~: 253\,000~USD avec un écart-type de 576\,000~USD et un équivalent certain de
585\,000~USD pour $\lambda=10^{-6}$. Combien coûte une erreur sur l'\omterm{def:mx:the-almgren-chriss-framework:urgency}{urgence}~? Un calendrier deux fois trop patient ($\kappa/2$) a un coût espéré de 258\,000~USD, avec un
écart-type de 542\,000, un équivalent certain de 551\,000, soit 5.8~\% de plus que l'optimum~; un calendrier deux fois trop pressé ($2\kappa$) a un coût espéré de 506\,000, avec un écart-type de 352\,000, un
équivalent certain de 630\,000, soit 20.9~\% de plus. La frontière est plate près de son optimum et raide du côté de la précipitation~: se tromper par excès de patience coûte moins cher.

\begin{omfigure}
\begin{tikzpicture}
\begin{groupplot}[group style={group size=2 by 1,horizontal sep=1.6cm},width=6.6cm,height=5cm,
  tick label style={font=\footnotesize},label style={font=\small},grid=major,grid style={gray!25},table/col sep=comma]
\nextgroupplot[title={\small part de l'ordre restant},xlabel={heures},xmin=0,xmax=6.5,ymin=0,ymax=1,
  legend columns=2,legend style={font=\footnotesize,at={(0.5,-0.3)},anchor=north,draw=none,fill=none}]
\addplot[black,thick] table[x=hours,y=k0]{figdata/microstructure/14-the-almgren-chriss-framework/trajectories.csv};
\addplot[omProp,thick] table[x=hours,y=k1]{figdata/microstructure/14-the-almgren-chriss-framework/trajectories.csv};
\addplot[omDef,very thick] table[x=hours,y=k2]{figdata/microstructure/14-the-almgren-chriss-framework/trajectories.csv};
\addplot[omThm,thick] table[x=hours,y=k3]{figdata/microstructure/14-the-almgren-chriss-framework/trajectories.csv};
\legend{$\kappa=0$,$\kappa=1$,$\kappa=1.99$,$\kappa=8$}
\nextgroupplot[title={\small frontière d'exécution efficiente (milliers de USD)},xlabel={écart-type},ylabel={coût espéré},xmin=0,xmax=620,ymin=150,ymax=700,
  legend columns=2,legend style={font=\footnotesize,at={(0.5,-0.3)},anchor=north,draw=none,fill=none}]
\addplot[omDef,very thick] table[x=sd,y=cost]{figdata/microstructure/14-the-almgren-chriss-framework/frontier.csv};
\addplot[omThm,only marks,mark=*] table[x=sd,y=cost]{figdata/microstructure/14-the-almgren-chriss-framework/points.csv};
\node[omThm,font=\scriptsize,anchor=west] at (axis cs:360,506) {pressé};
\node[omThm,font=\scriptsize,anchor=north east] at (axis cs:462,292) {optimal};
\node[omThm,font=\scriptsize,anchor=north] at (axis cs:535,240) {patient};
\node[omThm,font=\scriptsize,anchor=south] at (axis cs:585,262) {linéaire};
\end{groupplot}
\end{tikzpicture}

\omcaption{Vendre un million d'actions d'ici la clôture. À gauche~: trajectoires optimales pour plusieurs \omterm{def:mx:the-almgren-chriss-framework:urgency}{urgences} (par jour). À droite~: la \omterm{def:mx:the-almgren-chriss-framework:frontier}{frontière d'exécution efficiente}, avec
l'optimum pour $\lambda=10^{-6}$, la trajectoire pondérée par le temps, et des calendriers deux fois trop patient et deux fois trop pressé. Données~: \texttt{mx\_ac.block\_study},
\texttt{fig\_ac.py}.}\label{fig:mx:the-almgren-chriss-framework:frontier}
\end{omfigure}

\section{Calibrage et limites}

Le modèle comporte trois intrants, et chacun est une estimation. L'agent d'exécution de \texttt{mx\_ac} fait tourner un ordonnanceur dans le marché de \texttt{firm.agentmkt}~: il vend
20\,000 actions sur 30 minutes en 30 tranches, chacune jusqu'à la cible de l'ordonnanceur, pour des \omterm{def:mx:the-almgren-chriss-framework:urgency}{urgences} $\kappa T$ valant 0, 2 et 6, sur huit sessions chacune. Quant à
l'interface de l'ordonnanceur de \texttt{firm.acexec} (positions cibles à des instants donnés), c'est celle que les chapitres 15, 16 et 28 implémentent.

\omcode{code/microstructure/14-the-almgren-chriss-framework/python/mx_ac.py}{55}{71}{Un agent d'exécution~: à chaque intervalle, il vend jusqu'à la cible de l'ordonnanceur.}\label{lst:mx:the-almgren-chriss-framework:agent}

Le manque à gagner réalisé par action (\cref{fig:mx:the-almgren-chriss-framework:sim}) vaut 5.2 ticks pour la trajectoire linéaire (erreur type 2.8), 5.3 (2.3) pour
$\kappa T=2$ et 7.4 (1.3) pour $\kappa T=6$, avec des écarts types entre sessions de 8.0, 6.5 et 3.8~: les calendriers plus rapides coûtent davantage en moyenne et varient
moins, comme le prévoit le modèle. Calibré sur les mêmes essais, l'impact temporaire vaut $\eta=0.11$ tick par action et par action et par seconde (à partir du coût de chaque tranche
rapporté à sa taille) et $\sigma=0.69$ tick par racine de seconde (à partir des variations du milieu sur une minute), et le modèle prédit des moyennes de 1.2, 1.4 et 3.6 ticks et des
écarts types de 16.4, 13.2 et 7.6. Il ordonne correctement les calendriers, mais se trompe sur les niveaux de deux façons. La moyenne est trop basse, car le glissement de chaque tranche ne capte pas
l'impact qui s'accumule d'une tranche à l'autre~: l'écart implique un impact permanent d'environ $4\times10^{-4}$ ticks par action, soit le $\gamma$ mis à zéro ici. Quant à l'écart type
il est trop élevé, car le marché revient à la moyenne~: ses fondamentalistes ramènent le cours en arrière, si bien que la variance sur trente minutes est inférieure à trente
fois celle d'une minute. Almgren (2003) a étendu le modèle à un impact non linéaire et à un risque croissant avec le volume échangé~; le chapitre~15 laisse le calendrier réagir au
marché.

\begin{omfigure}
\begin{tikzpicture}
\begin{groupplot}[group style={group size=2 by 1,horizontal sep=1.6cm},width=6.6cm,height=5cm,
  tick label style={font=\footnotesize},label style={font=\small},grid=major,grid style={gray!25},table/col sep=comma,
  xlabel={urgence $\kappa T$},xtick={0,2,6},xmin=-0.5,xmax=6.5]
\nextgroupplot[title={\small manque à gagner moyen (ticks par action)},ymin=0,ymax=11,
  legend columns=1,legend cell align=left,legend style={font=\footnotesize,at={(0.5,-0.3)},anchor=north,draw=none,fill=none}]
\addplot[omDef,very thick,mark=*,error bars/.cd,y dir=both,y explicit] table[x=kt,y=mean,y error=se]{figdata/microstructure/14-the-almgren-chriss-framework/sim.csv};
\addplot[omThm,dashed,mark=square*] table[x=kt,y=model_mean]{figdata/microstructure/14-the-almgren-chriss-framework/sim.csv};
\legend{simulé,modèle}
\nextgroupplot[title={\small écart type (ticks par action)},ymin=0,ymax=18]
\addplot[omDef,very thick,mark=*] table[x=kt,y=sd]{figdata/microstructure/14-the-almgren-chriss-framework/sim.csv};
\addplot[omThm,dashed,mark=square*] table[x=kt,y=model_sd]{figdata/microstructure/14-the-almgren-chriss-framework/sim.csv};
\end{groupplot}
\end{tikzpicture}

\omcaption{Trajectoires d'Almgren--Chriss exécutées dans le marché simulé (huit sessions chacune), comparées au modèle calibré sur les mêmes essais. Le modèle
ordonne correctement les calendriers~; il ne tient pas compte de l'impact qui s'accumule d'une tranche à l'autre ni du retour à la moyenne du marché. Données~: \texttt{mx\_ac.sim\_study}.}\label{fig:mx:the-almgren-chriss-framework:sim}
\end{omfigure}

\section{Tutoriel~: liquider un bloc d'ici la clôture}\label{tut:mx:the-almgren-chriss-framework:tut}

\begin{tutorial}
\textbf{Objectif.} Résoudre le problème d'Almgren--Chriss pour un bloc, tracer sa frontière et tester les calendriers dans le marché simulé.
\textbf{État final~:} \cref{fig:mx:the-almgren-chriss-framework:frontier,fig:mx:the-almgren-chriss-framework:sim}.
\begin{tutsteps}
\item \textbf{Modèle.} \texttt{firm\_acexec.kappa}, \texttt{trajectory}, \texttt{discrete}, \texttt{cost\_var}, \texttt{frontier}, \texttt{time\_to}.
\item \textbf{Bloc.} \texttt{mx\_ac.block\_study(lam)}~: l'optimum, la trajectoire pondérée par le temps et les deux erreurs.
\item \textbf{Simulation.} \texttt{ScheduleAgent(ACScheduler(Q, 1, kT), \dots)} dans \texttt{firm\_agentmkt.session}~; \texttt{sim\_study()} pour trois
\omterm{def:mx:the-almgren-chriss-framework:urgency}{urgences} et huit graines, avec la calibration.
\item \textbf{Tracé.} \texttt{fig\_ac.py}.
\end{tutsteps}
\textbf{Pour aller plus loin.} Ajouter au modèle l'impact permanent impliqué par la simulation et recalculer ses moyennes~; calibrer $\sigma$ à
l'horizon de l'exécution au lieu d'une minute.
\end{tutorial}

\section{Construction~: le calendrier d'exécution}\label{bld:mx:the-almgren-chriss-framework:acexec}

\begin{build}
\bfield{Rôle} Le calendrier de référence de l'exécution de la société~: l'indice de référence que les chapitres 15 et 16 améliorent, l'interface de l'ordonnanceur que leurs
algorithmes et celui du chapitre~28 implémentent, ainsi que les estimations de coût et de risque avant négociation du chapitre~19.
\bfield{Interface} \texttt{kappa(eta, sigma, lam)}, \texttt{trajectory(X, T, kappa, t)}, \texttt{discrete(X, T, n, eta, sigma, lam)},
\texttt{cost\_var(holdings, T, gamma, eta, sigma)}, \texttt{frontier(\dots, lams)}, \texttt{time\_to(share, kappa, T)}~; \texttt{Scheduler.targets(times)},
\texttt{ACScheduler}, \texttt{LinearScheduler}.
\bfield{Règles} Les positions sont la quantité qui reste à négocier~; l'\omterm{def:mx:the-almgren-chriss-framework:urgency}{urgence} discrète à partir de $\cosh(\tilde\kappa\tau)$~; les coûts sont mesurés par rapport au prix d'arrivée~; $\eta$
par unité de taux de négociation dans l'unité de temps de l'appelant.
\bfield{Tests d'acceptation} \texttt{code/firm/acexec/tests/}~: les trajectoires discrète et continue concordent sur une grille fine~; le coût et la variance d'une trajectoire linéaire
calculés à la main~; les coûts de la frontière croissent et les variances décroissent avec $\lambda$~; le temps pour atteindre 95~\% vaut $\ln 20/\kappa$ sur un long horizon~; les cibles des ordonnanceurs.
\bfield{Pour aller plus loin} Impact temporaire non linéaire (Almgren, 2003)~; une version de portefeuille avec une matrice de covariance~; calibration à partir de \texttt{firm.impactfit}.
\end{build}

\begin{omsources}
\item D.~Bertsimas et A.~W.~Lo, ``Optimal control of execution costs'', \emph{Journal of Financial Markets} 1(1), 1998.
\item R.~Almgren et N.~Chriss, ``Optimal execution of portfolio transactions'', \emph{Journal of Risk} 3(2), 2001.
\item R.~Almgren, ``Optimal execution with nonlinear impact functions and trading-enhanced risk'', \emph{Applied Mathematical Finance} 10(1), 2003.
\end{omsources}

\section{Exercices}

\begin{exercise}[$\star$]\label{exo:mx:the-almgren-chriss-framework:1}
Avec $\eta=2.53\times10^{-7}$ dollar par action et par action et par jour, $\sigma=1$ dollar par action et par racine de jour, et $\lambda=10^{-6}$ par dollar, calculer $\kappa$.
\end{exercise}

\begin{exercise}[$\star$]\label{exo:mx:the-almgren-chriss-framework:2}
Calculer le coût espéré et l'écart type de la vente d'un million d'actions selon une trajectoire linéaire sur une journée, avec ces paramètres et sans impact permanent.
\end{exercise}

\begin{exercise}[$\star$]\label{exo:mx:the-almgren-chriss-framework:3}
Combien de temps le calendrier optimal met-il pour vendre 95~\% de l'ordre, et quelle est la formule lorsque la fenêtre est longue~?
\end{exercise}

\begin{exercise}[$\star\star$]\label{exo:mx:the-almgren-chriss-framework:4}
Pourquoi un calendrier deux fois trop pressé perd-il davantage d'équivalent certain (20.9~\%) qu'un calendrier deux fois trop patient (5.8~\%)~?
\end{exercise}

\begin{exercise}[$\star\star$]\label{exo:mx:the-almgren-chriss-framework:5}
Pourquoi le modèle calibré à partir du glissement par tranche prédit-il un manque à gagner moyen trop faible dans la simulation~?
\end{exercise}

\begin{exercise}[$\star\star$]\label{exo:mx:the-almgren-chriss-framework:6}
Pourquoi prédit-il un écart type trop élevé~?
\end{exercise}

\begin{exercise}[$\star\star\star$]\label{exo:mx:the-almgren-chriss-framework:7}
\emph{Programmation.} À partir des moyennes simulées et du modèle, déduire l'impact permanent $\gamma$ qui comblerait l'écart pour $\kappa T=0$ et pour $\kappa T=6$. Les
deux estimations sont-elles cohérentes~?
\end{exercise}

\begin{exercise}[$\star\star\star$]\label{exo:mx:the-almgren-chriss-framework:8}
\emph{Trouver la faille.} «~Un calendrier pondéré par le temps ne comporte aucun risque, car il négocie la même quantité chaque minute.~»
\end{exercise}

\section{Problème~: Liquider un bloc à la clôture}

\begin{problem}[{Problème du week-end --- liquider un bloc à la clôture}]\label{pb:mx:the-almgren-chriss-framework:1}
Un gérant de portefeuille doit vendre 1 million d'actions d'un titre à 50~dollars avant la clôture et demande au desk un calendrier, son coût espéré et son risque.

\medskip\noindent\textbf{Partie I --- Le modèle.}
\begin{enumerate}
\item Définir les \omterm{def:mx:the-almgren-chriss-framework:parent}{ordres parents}, les \omterm{def:mx:the-almgren-chriss-framework:parent}{ordres enfants} et une \omterm{def:mx:the-almgren-chriss-framework:trajectory}{trajectoire d'exécution}.
\item Écrire le coût espéré et la variance du \omterm{def:mx:the-almgren-chriss-framework:ac}{modèle d'Almgren--Chriss}.
\item Pourquoi le terme permanent ne dépend-il pas du calendrier~?
\item Dériver la trajectoire optimale.
\end{enumerate}

\medskip\noindent\textbf{Partie II --- Le bloc.}
\begin{enumerate}[resume]
\item Comment $\eta$ est-il linéarisé à partir de la loi en racine carrée, et que vaut-il ici~?
\item Donner $\kappa$, le coût espéré, l'écart type et l'équivalent certain pour $\lambda=10^{-6}$.
\item Quand 95~\% de l'ordre est-il vendu~?
\item Donner le coût, le risque et l'équivalent certain de la trajectoire linéaire.
\end{enumerate}

\medskip\noindent\textbf{Partie III --- La frontière.}
\begin{enumerate}[resume]
\item Définir la \omterm{def:mx:the-almgren-chriss-framework:frontier}{frontière d'exécution efficiente}.
\item Donner le coût, le risque et l'équivalent certain des calendriers deux fois trop patients et deux fois trop pressés.
\item Quelle erreur coûte le moins cher, et pourquoi~?
\item Que signifie l'aversion au risque pour le gérant~?
\end{enumerate}

\medskip\noindent\textbf{Partie IV --- Calibration et verdict.}
\begin{enumerate}[resume]
\item Décrire le test simulé.
\item Donner les moyennes et les écarts types réalisés.
\item Donner les paramètres calibrés du modèle et ses prédictions.
\item Expliquer les écarts.
\item Quel impact permanent comble l'écart de moyenne~?
\item Énoncer le \emph{résultat nommé}~: l'horizon d'exécution optimal, le coût espéré et l'écart type pour une aversion au risque donnée, ainsi que le surcoût d'un calendrier
deux fois trop patient ou deux fois trop pressé.
\item Que diriez-vous au gérant sur la fiabilité des chiffres~?
\item En une phrase~: qu'apporte la frontière à une estimation de coût~?
\end{enumerate}
\end{problem}

\section{Questions d'entretien}

\begin{interviewq}[$\star$ \iqroles{trader, chercheur}]\label{iq:mx:the-almgren-chriss-framework:1}
Expliquez le compromis de l'exécution optimale et la forme de la trajectoire d'Almgren--Chriss.
\end{interviewq}

\begin{interviewq}[$\star\star$ \iqroles{chercheur}]\label{iq:mx:the-almgren-chriss-framework:2}
Dérivez la trajectoire d'Almgren--Chriss à partir de son objectif.
\end{interviewq}

\begin{interviewq}[$\star\star$ \iqroles{trader}]\label{iq:mx:the-almgren-chriss-framework:3}
Votre client double son aversion au risque. Que se passe-t-il pour le calendrier, le coût espéré et le risque~?
\end{interviewq}

\begin{interviewq}[$\star\star$ \iqroles{chercheur}]\label{iq:mx:the-almgren-chriss-framework:4}
Comment calibreriez-vous l'impact temporaire et l'impact permanent du modèle~?
\end{interviewq}

\begin{interviewq}[$\star\star$ \iqroles{développeur}]\label{iq:mx:the-almgren-chriss-framework:5}
Concevez une interface de planificateur que différents algorithmes d'exécution peuvent partager. Que prend-elle en entrée et que renvoie-t-elle~?
\end{interviewq}

\begin{interviewq}[$\star\star\star$ \iqroles{chercheur}]\label{iq:mx:the-almgren-chriss-framework:6}
Que laisse de côté le \omterm{def:mx:the-almgren-chriss-framework:ac}{modèle d'Almgren--Chriss}, et lequel de ces éléments compte le plus pour une action liquide~?
\end{interviewq}
